\chapter{Abelian varieties}\label{ch:abelian}

Abelian varieties are the one large class of varieties whose Hodge classes
can be written down. The reason is that their cohomology is an exterior
algebra on the first cohomology, so that everything is decided by a single
Hodge structure of weight one and the group that preserves it. This
chapter sets up that dictionary. The references for the background are
Mumford's book \cite{Mum70} and the book of Birkenhake and Lange
\cite{BL04}. We prove the facts about Hodge classes, since they are the
point of the chapter.

\section{Complex tori and their cohomology}

Let $\mathfrak{t}$ be a complex vector space of dimension $g$ and let $\Lambda\subset
\mathfrak{t}$ be a lattice, that is, a discrete subgroup of rank $2g$. The quotient
$T=\mathfrak{t}/\Lambda$ is a compact complex manifold and a commutative complex Lie
group, and it is called a \emph{complex torus}. Its fundamental group is
$\Lambda$, and as a real manifold it is a product of $2g$ circles.

\begin{proposition}\label{prop:torcoh}
Cup product induces isomorphisms
\[
  \textstyle\bigwedge^{k}H^{1}(T,\ZZ)\xrightarrow{\;\sim\;}H^{k}(T,\ZZ),
  \qquad H^{1}(T,\ZZ)=\Hom(\Lambda,\ZZ).
\]
With complex coefficients, $H^{1,0}(T)$ is the space of $\CC$-linear forms
on $\mathfrak{t}$, $H^{0,1}(T)$ is the space of $\CC$-antilinear forms, and
\[
  H^{p,q}(T)=\textstyle\bigwedge^{p}H^{1,0}(T)\otimes\bigwedge^{q}H^{0,1}(T).
\]
\end{proposition}

\begin{proof}
For a product of circles, the K\"unneth formula shows that the integral
cohomology ring is the exterior algebra on the degree-one classes dual to
the circles. These classes are a basis of $\Hom(\Lambda,\ZZ)$. For the
second part, give $T$ the flat metric coming from a hermitian inner
product on $\mathfrak{t}$. A form with constant coefficients in the linear coordinates
$z_{1},\dots,z_{g}$ is harmonic, and the constant forms
$dz_{I}\wedge d\bar z_{J}$ with $|I|=p$ and $|J|=q$ span a space of
dimension $\binom{g}{p}\binom{g}{q}$. Summing over $p+q=k$ gives
$\binom{2g}{k}$, which is the $k$-th Betti number, so these forms are all
the harmonic forms, and their types give the Hodge decomposition. In degree
one, $dz_{j}$ is a $\CC$-linear form and $d\bar z_{j}$ an antilinear one.
\end{proof}

So the Hodge structure of $T$ in every degree is determined by the one in
degree one. Conversely, an integral Hodge structure of weight one, of
type $\{(1,0),(0,1)\}$, determines a complex torus: one takes the dual of
$H^{1,0}$ modulo the lattice $H_{1}$. Complex tori and integral Hodge
structures of this type are the same thing, and homomorphisms of tori are
the same as morphisms of these Hodge structures.

\section{Line bundles and polarisations}

By Proposition~\ref{prop:torcoh}, $H^{2}(T,\ZZ)=\bigwedge^{2}\Hom(\Lambda,\ZZ)$
is the group of alternating forms $E\colon\Lambda\times\Lambda\to\ZZ$. Write
$E_{\RR}$ for the extension of $E$ to $\mathfrak{t}$, viewed as a real vector space,
and $i$ for the complex structure of $\mathfrak{t}$.

\begin{proposition}\label{prop:ns}
An alternating form $E$ on $\Lambda$ is of type $(1,1)$, as a class in
$H^{2}(T,\CC)$, if and only if $E_{\RR}(ix,iy)=E_{\RR}(x,y)$ for all
$x,y\in \mathfrak{t}$. Every such integral form is the first Chern class of a
holomorphic line bundle on $T$.
\end{proposition}

\begin{proof}
Extend $E$ complex-bilinearly to $\mathfrak{t}\otimes_{\RR}\CC=\mathfrak{t}'\oplus \mathfrak{t}''$, the
$(+i)$ and $(-i)$ eigenspaces of the complex structure. A class in
$\bigwedge^{2}$ of the dual has type $(1,1)$ exactly when it vanishes on
$\mathfrak{t}'\times \mathfrak{t}'$ and on $\mathfrak{t}''\times \mathfrak{t}''$. If $E$ is invariant under $i$, then for
$x,y\in \mathfrak{t}'$ we get $E(x,y)=E(ix,iy)=i^{2}E(x,y)$, so $E(x,y)=0$, and the
same on $\mathfrak{t}''$. Conversely, if $E$ vanishes on both, write $x=x'+x''$ and
$y=y'+y''$. Then
\[
  E(ix,iy)=E(ix',-iy'')+E(-ix'',iy')=E(x',y'')+E(x'',y')=E(x,y).
\]
The second statement is Steps~1 and~2 of the proof of the Lefschetz
theorem in Chapter~\ref{ch:lefschetz}, which use only that $T$ is a
compact K\"ahler manifold.
\end{proof}

To such a form $E$ attach the function
\[
  H(x,y)=E_{\RR}(ix,y)+iE_{\RR}(x,y),\qquad x,y\in \mathfrak{t}.
\]
A direct check using the invariance of $E_{\RR}$ under $i$ shows that $H$
is a hermitian form, $\CC$-linear in the first variable, with imaginary
part $E$. The N\'eron--Severi group of $T$ is the group of these forms.

\begin{theorem}[Riemann, Lefschetz]\label{thm:riemann}
The torus $T$ is projective if and only if there is an integral
alternating form $E$, invariant under $i$, whose hermitian form $H$ is
positive definite. In that case $E$ is the Chern class of an ample line
bundle.
\end{theorem}

One direction is easy. If $T$ embeds in projective space, the restriction
of the Fubini--Study form is a positive $(1,1)$ form, its average under
translations has constant coefficients and the same class, and its
hermitian form is positive definite. The other direction is the theorem
of Lefschetz that the third power of such a line bundle is very ample. It
is proved with theta functions \cite{Mum70}.

An \emph{abelian variety} is a projective complex torus. A
\emph{polarisation} is the class of an ample line bundle, or equivalently
a form $E$ as in Theorem~\ref{thm:riemann}. The form $E$ is a
polarisation, in the sense of Chapter~\ref{ch:toolkit}, of the Hodge
structure $H_{1}(A,\QQ)=\Lambda_{\QQ}$. Since $E$ is nondegenerate it
identifies $\Lambda_{\QQ}$ with its dual $H^{1}(A,\QQ)$, and we use the
same letter for the induced form on $H^{1}(A,\QQ)$. From now on we write
\[
  V=H^{1}(A,\QQ),
\]
so that $H^{k}(A,\QQ)=\bigwedge^{k}V$.

\section{Endomorphisms and the Rosati involution}

A homomorphism of abelian varieties is a holomorphic map that is a group
homomorphism. It lifts to a $\CC$-linear map of universal covers that
carries one lattice into the other. Write $\End(A)$ for the ring of
endomorphisms and $\End^{0}(A)=\End(A)\otimes\QQ$. An element of
$\End^{0}(A)$ acts on $\Lambda_{\QQ}$, and its real extension commutes with
$i$. Conversely, a $\QQ$-linear map of $\Lambda_{\QQ}$ whose real extension
commutes with $i$ has an integer multiple preserving $\Lambda$, so it lies
in $\End^{0}(A)$. In other words, $\End^{0}(A)$ is the algebra of
endomorphisms of the Hodge structure $H_{1}(A,\QQ)$.

By Poincar\'e's complete reducibility theorem, every abelian variety is
isogenous to a product of powers of simple, pairwise non-isogenous abelian
varieties, and $\End^{0}(A)$ is a semisimple $\QQ$-algebra
\cite{Mum70}.

Fix a polarisation $E$. For $\varphi\in\End^{0}(A)$, acting on
$\Lambda_{\QQ}$, let $\varphi'$ be its adjoint with respect to $E$:
\[
  E(\varphi x,y)=E(x,\varphi'y).
\]

\begin{proposition}\label{prop:rosati}
The adjoint $\varphi'$ lies in $\End^{0}(A)$, and the map
$\varphi\mapsto\varphi'$ is an involution of $\End^{0}(A)$, called the
\emph{Rosati involution}. It is positive: $\Tr(\varphi\varphi')>0$ for
$\varphi\ne0$, the trace being taken on $\Lambda_{\QQ}$.
\end{proposition}

\begin{proof}
Invariance of $E$ under $i$ gives $E(iu,v)=-E(u,iv)$ for all $u,v$. Using
this twice, and that $\varphi$ commutes with $i$,
\[
  E(x,\varphi'iy)=E(\varphi x,iy)=-E(i\varphi x,y)=-E(\varphi ix,y)
  =-E(ix,\varphi'y)=E(x,i\varphi'y),
\]
so $\varphi'$ commutes with $i$ and lies in $\End^{0}(A)$. It is clearly an
involution. For positivity, a similar computation shows
$H(\varphi x,y)=H(x,\varphi'y)$, so $\varphi'$ is the adjoint of $\varphi$
for the positive definite hermitian form $H$. The trace of $\varphi\varphi'$
on the real space $\mathfrak{t}$ is twice the real part of its complex trace, and the
complex trace of $\varphi\varphi^{*}$ is the sum of the squared lengths of
the columns of $\varphi$ in an $H$-orthonormal basis, which is positive.
\end{proof}

\begin{proposition}\label{prop:nssym}
The map that sends a symmetric element $\varphi=\varphi'$ of $\End^{0}(A)$
to the form $E_{\varphi}(x,y)=E(\varphi x,y)$ is an isomorphism
\[
  \{\varphi\in\End^{0}(A):\varphi'=\varphi\}\xrightarrow{\;\sim\;}
  \NS(A)\otimes\QQ .
\]
\end{proposition}

\begin{proof}
If $\varphi$ is symmetric, $E_{\varphi}$ is alternating, because
$E(\varphi y,x)=E(y,\varphi x)=-E(\varphi x,y)$. It is invariant under $i$,
because $\varphi$ commutes with $i$ and $E$ is invariant. So it lies in
$\NS(A)_{\QQ}$ by Proposition~\ref{prop:ns}. Conversely, given an
alternating form $E'$ that is rational and invariant under $i$, the
nondegeneracy of $E$ gives a unique $\QQ$-linear $\varphi$ with
$E'(x,y)=E(\varphi x,y)$. The computation in the proof of
Proposition~\ref{prop:rosati}, applied to $E'$ and $E$, shows that $\varphi$
commutes with $i$, so it lies in $\End^{0}(A)$, and it is symmetric because
$E'$ is alternating.
\end{proof}

So the divisor classes of $A$ are controlled by its endomorphisms. This is
the first sign of a principle that runs through the book: whatever is
built from divisors is invisible to anything that commutes with the
endomorphisms and preserves the polarisation.

The simple case is governed by Albert's classification. If $A$ is simple,
$\End^{0}(A)$ is a division algebra of finite dimension over $\QQ$ with a
positive involution, and such algebras are of four types: a totally real
field; a totally indefinite or a totally definite quaternion algebra over
a totally real field; or a division algebra over a CM field, with an
involution that restricts to complex conjugation on the centre
\cite{Mum70}. Abelian varieties of Weil type, the subject of
Chapter~\ref{ch:weil}, have an endomorphism algebra of the fourth type in
its simplest form, an imaginary quadratic field.

\section{The Hodge group of an abelian variety}

Write $\Hg(A)$ for the Hodge group of the Hodge structure $V=H^{1}(A,\QQ)$,
as defined in Section~\ref{sec:mt}.

\begin{proposition}\label{prop:hgabelian}
\begin{enumerate}[label=(\roman*)]
\item $\Hg(A)$ is contained in the symplectic group $\mathrm{Sp}(V,E)$.
\item The commutant of $\Hg(A)$ in $\End_{\QQ}(V)$ is $\End^{0}(A)$, acting on
$V$ by pullback.
\item The Hodge classes on $A$ are the invariants of $\Hg(A)$:
\[
  \Hdg^{*}(A)=\Bigl(\textstyle\bigwedge^{*}V\Bigr)^{\Hg(A)}.
\]
\item For every $m\ge1$, the Hodge group of $A^{m}$ is $\Hg(A)$ acting
diagonally on $H^{1}(A^{m},\QQ)=V^{\oplus m}$.
\end{enumerate}
\end{proposition}

\begin{proof}
(i) The polarisation is a morphism of Hodge structures
$V\otimes V\to\QQ(-1)$. On $\QQ(-1)$, of type $(1,1)$, an element $z$ of
the unit circle acts by $z\bar z=1$. So $h(z)$ preserves $E$ for $|z|=1$.
The stabiliser of $E$ is the algebraic group $\mathrm{Sp}(V,E)$, defined over
$\QQ$, so it contains $\Hg(A)$ by minimality.

(ii) The centraliser of a $\QQ$-linear map $\varphi$ is an algebraic
subgroup defined over $\QQ$, so $\varphi$ commutes with $\Hg(A)$ if and only
if it commutes with $h(z)$ for all $|z|=1$. The element $h(z)$ acts on
$V^{1,0}$ and $V^{0,1}$ by different scalars when $z\ne\pm1$, so
commuting with all of them means preserving the Hodge decomposition. That
is, $\varphi$ is an endomorphism of the Hodge structure $V$, which by the
dictionary of the last section means an element of $\End^{0}(A)$.

(iii) The space $\bigwedge^{k}V$ is the image of the antisymmetrisation
projector on $V^{\otimes k}$, which is both a morphism of Hodge structures
and $\GL(V)$-equivariant. So an element of $\bigwedge^{k}V$ is a Hodge class,
or an invariant, if and only if it is one in $V^{\otimes k}$. Now apply
Theorem~\ref{thm:hginvariants}.

(iv) The homomorphism defining the Hodge structure of $A^{m}$ is the
diagonal $h\oplus\dots\oplus h$. The diagonal embedding
$\GL(V)\to\GL(V^{\oplus m})$ is a closed embedding defined over $\QQ$, so the
smallest $\QQ$-algebraic subgroup containing the diagonal image of $h(U(1))$
is the diagonal image of $\Hg(A)$.
\end{proof}

Part (iv) means that the representation of one reductive group on one
vector space decides the Hodge classes not only on $A$ but on all its
powers, and hence on every abelian variety isogenous to a product of
copies of $A$. Part (iii) turns the computation of Hodge classes into
invariant theory.

For a very general principally polarised abelian variety the Hodge group
is the whole symplectic group. One way to see this is through monodromy.
Over the moduli space with level structure, the monodromy of $H^{1}$ is a
subgroup of finite index in $\mathrm{Sp}_{2g}(\ZZ)$, which is Zariski dense in
$\mathrm{Sp}_{2g}$. By Andr\'e's theorem \cite{And92}, the algebraic monodromy
group is contained in the Hodge group of a very general member. So that
Hodge group contains $\mathrm{Sp}_{2g}$, and by part~(i) it equals it.

\section{The Lefschetz group and Lefschetz classes}

\begin{definition}
The \emph{Lefschetz group} $L(A)$ is the connected component of the
identity of the centraliser of $\End^{0}(A)$ in $\mathrm{Sp}(V,E)$. The
\emph{Lefschetz classes} of $A$ are the elements of the subalgebra
$\mathcal D^{*}(A)\subseteq H^{*}(A,\QQ)$ generated by divisor classes.
\end{definition}

By Proposition~\ref{prop:hgabelian}(i) and (ii), and because $\Hg(A)$ is
connected,
\[
  \Hg(A)\subseteq L(A).
\]
The Lefschetz group is easy to compute, since it depends only on the
endomorphism algebra with its involution. The Hodge group is hard to
compute, since it depends on the period point.

\begin{proposition}\label{prop:lefinv}
Every divisor class on $A$ is invariant under $L(A)$. Hence
\[
  \mathcal D^{*}(A)\subseteq\Bigl(\textstyle\bigwedge^{*}V\Bigr)^{L(A)}
  \subseteq\Bigl(\textstyle\bigwedge^{*}V\Bigr)^{\Hg(A)}=\Hdg^{*}(A).
\]
\end{proposition}

\begin{proof}
It is enough to work on $H_{1}$, where $L(A)$ acts through the
contragredient. By Proposition~\ref{prop:nssym}, a divisor class is a form
$E(\varphi x,y)$ with $\varphi$ in $\End^{0}(A)$. An element $g$ of $L(A)$
preserves $E$ and commutes with $\varphi$, so
$E(\varphi gx,gy)=E(g\varphi x,gy)=E(\varphi x,y)$. The middle inclusion
holds because $\Hg(A)\subseteq L(A)$, and the last equality is
Proposition~\ref{prop:hgabelian}(iii).
\end{proof}

The first inclusion is in fact an equality. This theorem is due to Murty
and Milne \cite{Mur84,Mil99}: the Lefschetz classes are exactly the
invariants of the Lefschetz group. We shall use only the inclusion proved
above, together with the following consequence of the theorem.

\begin{corollary}\label{cor:hgeqlef}
If $\Hg(A)$ and $L(A)$ have the same invariants in $\bigwedge^{*}V$, and in
particular if $\Hg(A)=L(A)$, then every Hodge class on $A$ is a polynomial
in divisor classes, and the Hodge conjecture holds for $A$.
\end{corollary}

This covers a great deal. For a very general principally polarised
abelian variety, both groups are the symplectic group, and the Hodge ring
is generated by the class of the theta divisor. Many classes of abelian
varieties are known to satisfy the hypothesis of the corollary, among them
simple abelian varieties of prime dimension and products of elliptic
curves \cite{Rib83,BL04}. Moonen and Zarhin worked through all abelian
varieties of dimension at most five, and in that range the Hodge classes
that are not Lefschetz classes are all accounted for by Weil classes
\cite{MZ99}.

The first place where the hypothesis fails is the subject of the next
chapter. There the endomorphism algebra is an imaginary quadratic field
$K$, the Lefschetz group is a unitary group, and the Hodge group is the
special unitary group, one dimension smaller. The difference is a single
character, the determinant, and it is exactly enough to create Hodge
classes that no divisor can explain.
